\documentclass[11pt]{article}
\usepackage[utf8]{inputenc}
\usepackage[T1]{fontenc}
\usepackage{parskip}
\usepackage[margin=1in]{geometry}

\begin{document}

\noindent Eike Natan Sousa Brito \\
Universidade Federal de Sergipe \\
Cidade Universit\'aria Prof. Jos\'e Alo\'isio de Campos \\
S\~ao Crist\'ov\~ao, Sergipe, 49100-000, Brazil \\
eike.sousa@hotmail.com

\bigskip
\noindent June 16, 2026

\bigskip
\noindent To the Editors \\
\textit{MethodsX}

\bigskip
\noindent Dear Editors,

We are pleased to submit our manuscript entitled ``A Comparative Analysis of
Open-Source LLMs for Converting Architecture, Engineering and Construction
Standards to the RASE Format'' for consideration in \textit{MethodsX}.

In line with the scope of \textit{MethodsX}, this manuscript describes a
reproducible method rather than a single set of results. We present RASE-PE, a
three-level normalization pipeline that uses Prompt Engineering alone---without
fine-tuning---to convert Architecture, Engineering and Construction (AEC)
technical standards into the computable RASE format (Requirement, Applicability,
Selection, Exception) required for automated BIM compliance checking. The method
decomposes the conversion into three successive levels: (N1) atomic segmentation
of normative sentences, (N2) identification of RASE operators, and (N3) semantic
JSON structuring.

We believe the method fits \textit{MethodsX} for the following reasons:

\begin{itemize}
  \item It is a transferable, step-by-step procedure that other researchers and
        practitioners can reproduce and adapt to their own standards and models.
  \item It was validated on a real dataset of 79 annotated NBR~9050 standards
        across five experiments, using thirteen metrics from five families
        (lexical, semantic, distance-based, generation-oriented, and
        classification), reaching semantic similarity close to 0.90 in the best
        scenarios.
  \item It relies exclusively on open-source LLMs (Llama, Dolphin, Gemma,
        Mistral, Alpaca, and Qwen) served locally through Ollama, under an
        identical dataset, prompt, and hardware protocol, so the full procedure
        can be reproduced without proprietary tools.
\end{itemize}

The manuscript follows the \textit{MethodsX} article structure, includes a
CRediT author statement and a data availability statement, and reports the
method in sufficient detail to be repeated.

We confirm that this work is original, has not been published previously, and is
not under consideration by any other journal. All authors have approved the
submission. The authors declare that they have no known competing financial
interests or personal relationships that could have influenced the work reported
in this paper.

Thank you for considering our submission.

\bigskip
\noindent Sincerely,

\bigskip
\noindent Eike Natan Sousa Brito (on behalf of all authors) \\
Co-authors: Andr\'e Carvalho; Marco Ant\^onio Brasiel Sampaio \\
Universidade Federal de Sergipe, S\~ao Crist\'ov\~ao, Sergipe, Brazil \\
eike.sousa@hotmail.com

\end{document}
